% WWW 2027 Research Track paper template.
% Official review requirement:
%   \documentclass[sigconf, anonymous, review]{acmart}
\documentclass[sigconf,anonymous,review]{acmart}

% ---------------------------------------------------------------------------
% Submission mode and packages
% ---------------------------------------------------------------------------
% Keep page numbers in the review PDF. Do not change the paper to one-column
% or manuscript mode: WWW 2027 requires the ACM double-column format.
\settopmatter{printfolios=true}

% Mathematics
\usepackage{amsmath}
\usepackage{mathtools}
% acmart already loads the AMS symbol fonts; do not reload amssymb here.

% Figures and tables
\usepackage{graphicx}
\usepackage{booktabs}
\usepackage{multirow}
\usepackage{tabularx}
\usepackage{array}
\usepackage{makecell}
\usepackage{xcolor}

% Algorithms and code listings
\usepackage{algorithm}
\usepackage{algpseudocode}
\usepackage{listings}

% Balance the final page when the paper is ready.
\usepackage{balance}

% The acmart class already loads hyperref, url, microtype, caption, xcolor,
% graphicx, natbib, and several supporting packages. Do not reload them with
% custom options unless the ACM template documentation explicitly requires it.

\lstset{
  basicstyle=\ttfamily\footnotesize,
  columns=fullflexible,
  frame=single,
  breaklines=true,
  showstringspaces=false,
  captionpos=b
}

% ---------------------------------------------------------------------------
% Title and author information
% ---------------------------------------------------------------------------
% Do not reveal authors, affiliations, acknowledgements, funding, or other
% identifying information in the anonymous review version.
\title[Short Paper Title]{Full Paper Title}
\author{Anonymous Author(s)}
\renewcommand{\shortauthors}{Anonymous Author(s)}

% Camera-ready metadata only. Keep these commented during anonymous review.
% \acmConference[The Web Conference '27]
%   {The Web Conference 2027}{May 10--14, 2027}{Dublin, Ireland}
% \copyrightyear{2027}
% \acmYear{2027}
% \setcopyright{acmlicensed}

\begin{document}

% ---------------------------------------------------------------------------
% Abstract, CCS concepts, and keywords
% ---------------------------------------------------------------------------
\begin{abstract}
State the problem, the proposed approach, the main evidence, and the central
result in a self-contained abstract. Do not use a placeholder or dummy
abstract in the submitted version.
\end{abstract}

% Replace these examples with the CCS concepts that match the paper.
\ccsdesc[500]{Information systems~Web mining}
\ccsdesc[300]{Computing methodologies~Machine learning}
\keywords{Web research, machine learning, reproducibility}

\maketitle

% ---------------------------------------------------------------------------
% 1. Introduction
% ---------------------------------------------------------------------------
\section{Introduction}
\label{sec:introduction}

Explain the Web-related problem, why it matters, and the gap in existing
research. WWW 2027 requires the first page to state clearly how the work is
relevant to the Web and to the selected track.

\textbf{Contributions.} Summarize the contributions in a short list:
\begin{itemize}
  \item Contribution one: state the main scientific or technical advance.
  \item Contribution two: state the method, system, dataset, or analysis.
  \item Contribution three: state the empirical, theoretical, or practical
        evidence supporting the claims.
\end{itemize}

% ---------------------------------------------------------------------------
% 2. Related work
% ---------------------------------------------------------------------------
\section{Related Work}
\label{sec:related-work}

Position the work against the closest Web, machine learning, and application
literature. Explain precisely what is different from prior work rather than
only listing papers.

% ---------------------------------------------------------------------------
% 3. Web relevance and problem formulation
% ---------------------------------------------------------------------------
\section{Web Relevance and Problem Formulation}
\label{sec:problem}

\subsection{Web Relevance}
\label{sec:web-relevance}

State the Web object, platform, infrastructure, socio-economic system, or
Web-related impact studied by the paper. Make this connection explicit and
concrete.

\subsection{Problem Definition}
\label{sec:problem-definition}

Define the task, inputs, outputs, assumptions, notation, and evaluation
setting. For example, a generic objective can be written as
\begin{equation}
  \mathcal{L}(\theta)
  =
  \mathcal{L}_{\mathrm{task}}(\theta)
  + \lambda \mathcal{R}(\theta),
  \label{eq:objective}
\end{equation}
where $\theta$ denotes the learnable parameters and $\lambda$ controls the
regularization strength. Replace this example with the paper's actual
formulation.

% ---------------------------------------------------------------------------
% 4. Method
% ---------------------------------------------------------------------------
\section{Method}
\label{sec:method}

\subsection{Overview}
\label{sec:method-overview}

Describe the full pipeline and explain how each component addresses the
problem defined in Section~\ref{sec:problem}.

\begin{figure}[t]
  \centering
  % Replace the box with a real figure, for example:
  % \includegraphics[width=\columnwidth]{figures/method-overview.pdf}
  \fbox{\rule{0pt}{1.15in}\rule{0.95\columnwidth}{0pt}}
  \caption{Placeholder for the main method, system, or data-flow overview.}
  \Description{A placeholder diagram for the main method, system, or data-flow overview.}
  \label{fig:method-overview}
\end{figure}

\subsection{Technical Details}
\label{sec:technical-details}

Give enough detail to reproduce the method, including data processing,
training or optimization, hyperparameters, computational resources, and
implementation choices.

\begin{algorithm}[t]
  \caption{Generic research procedure}
  \label{alg:procedure}
  \begin{algorithmic}[1]
    \Require Input data $\mathcal{D}$ and configuration $\lambda$
    \Ensure Trained model or analysis result $\hat{y}$
    \State Initialize $\theta$
    \For{$t \gets 1$ to $T$}
      \State Compute the objective $\mathcal{L}(\theta;\mathcal{D},\lambda)$
      \State Update $\theta$
    \EndFor
    \State \Return $\hat{y}$
  \end{algorithmic}
\end{algorithm}

% ---------------------------------------------------------------------------
% 5. Experimental setup
% ---------------------------------------------------------------------------
\section{Experimental Setup}
\label{sec:experiments}

\subsection{Datasets and Tasks}
\label{sec:datasets}

Describe the data source, collection period, preprocessing, splits, licenses,
and any filtering or deduplication. Report ethical review or consent
considerations when human data are involved.

\subsection{Baselines and Metrics}
\label{sec:baselines}

Define every baseline, implementation source, tuning protocol, metric, and
aggregation rule. State whether higher or lower values are better.

\subsection{Implementation Details}
\label{sec:implementation}

Report hardware, software versions, random seeds, training budget, stopping
criteria, and statistical testing or uncertainty estimation.

\begin{table}[t]
  \caption{Placeholder for the main quantitative results.}
  \label{tab:results}
  \centering
  \begin{tabular}{lcc}
    \toprule
    Method & Metric 1 $\uparrow$ & Metric 2 $\downarrow$ \\
    \midrule
    Baseline & 0.00 & 0.00 \\
    Proposed method & 0.00 & 0.00 \\
    \bottomrule
  \end{tabular}
\end{table}

\subsection{Reproducibility}
\label{sec:reproducibility}

State what code, data, configuration, and artifacts will be released, and
where reviewers can access any artifact required to verify the claims.

% ---------------------------------------------------------------------------
% 6. Results
% ---------------------------------------------------------------------------
\section{Results}
\label{sec:results}

Present the main results first, followed by uncertainty or statistical tests.
Use precise comparisons and tie each result to a stated research question or
hypothesis. Refer to Figure~\ref{fig:method-overview} and
Table~\ref{tab:results} using labels rather than hard-coded numbers.

\begin{lstlisting}[language=Python,caption={Placeholder for a short reproducibility-critical code fragment.},label={lst:code}]
def evaluate(model, dataset):
    predictions = model.predict(dataset)
    return summarize_metrics(predictions)
\end{lstlisting}

% ---------------------------------------------------------------------------
% 7. Analysis
% ---------------------------------------------------------------------------
\section{Analysis}
\label{sec:analysis}

\subsection{Ablation and Sensitivity}
\label{sec:ablation}

Remove or vary important components and explain which conclusions remain
stable.

\subsection{Qualitative or Error Analysis}
\label{sec:error-analysis}

Inspect representative successes and failures, identify limitations, and
report any subgroup or distributional differences that affect the claims.

% ---------------------------------------------------------------------------
% 8. Limitations and ethics
% ---------------------------------------------------------------------------
\section{Limitations and Ethical Considerations}
\label{sec:limitations}

State the scope conditions, threats to validity, known failure modes, data
limitations, privacy and consent issues, misuse risks, and any environmental
or societal impact relevant to the work.

% ---------------------------------------------------------------------------
% 9. Conclusion
% ---------------------------------------------------------------------------
\section{Conclusion}
\label{sec:conclusion}

Summarize the problem, the main contribution, the strongest evidence, and the
remaining open questions without introducing new claims.

% ---------------------------------------------------------------------------
% References
% ---------------------------------------------------------------------------
% Use the ACM reference style. Create references.bib and uncomment these two
% lines when the bibliography is ready:
% \bibliographystyle{ACM-Reference-Format}
% \bibliography{references}

% Balance the final page after the main text is complete. If an appendix is
% substantial, move this command to the final page of the main paper.
% \balance

% ---------------------------------------------------------------------------
% Optional appendix
% ---------------------------------------------------------------------------
\appendix
\section{Additional Details}
\label{app:details}

Use the optional appendix for proofs, pseudocode, extended experimental
details, or reproducibility information. The first eight pages of a long
paper should remain self-contained because reviewers are not required to read
past the main-paper limit.

\subsection{Extended Results}
\label{app:extended-results}

Add supplementary tables, derivations, or additional analyses here when they
are necessary to support the paper's claims.

\end{document}

% ---------------------------------------------------------------------------
% WWW 2027 anonymous-submission checklist
% ---------------------------------------------------------------------------
% [ ] English, ACM double-column format, and a single PDF.
% [ ] Long paper: 8 pages of main paper; at most 12 pages including references
%     and an optional appendix. Short paper: 4 pages total.
% [ ] The first page explicitly explains relevance to the Web and the track.
% [ ] No author names, affiliations, acknowledgements, funding, self-identifying
%     URLs, or other identifying information appear in the review version.
% [ ] Self-citations are written in the third person where needed for anonymity.
% [ ] No dummy abstract; all authors, conflicts, and the OpenReview profile are
%     complete before submission.
% [ ] Check the current WWW 2027 deadline page before submitting.
